problem:  
Let \( (M, \Delta, \|\cdot\|_x) \) be a connected, equiregular sub‑Finsler manifold of step \(s\ge2\), with smooth bracket‑generating distribution \( \Delta\subset TM \) and measurable, strictly convex horizontal norms \(x\mapsto \|\cdot\|_x\) on \(\Delta_x\).  
For each \(x\in M\), let the metric tangent cone \(\mathrm{Tan}(M,d_\varphi,x)\) be a Carnot group \(G_x\) with graded Lie algebra \(\mathfrak g_x=V^x_1\oplus\cdots\oplus V^x_s\) and left‑invariant norm induced by \(\|\cdot\|_x\) on \(V^x_1\).

Assume that for \(\mathfrak m\)-almost every \(x\), the tangent Carnot group \(G_x\) has graded Lie algebra **isomorphic to a fixed model** \(\mathfrak g=V_1\oplus\cdots\oplus V_s\) (so their bracket structures agree up to linear isomorphism), but **the horizontal norm \(\|\cdot\|_x\)** may vary measurably over \(V_1\).  

**Question:**  
Can one construct, or preclude, a smooth equiregular sub‑Finsler manifold for which:  
- all graded Lie algebras of the tangent cones are isomorphic to \(\mathfrak g\) almost everywhere,  
- the metric tangents themselves all take the form of Carnot groups modeled on \(\mathfrak g\), but  
- the manifold fails to be locally bi‑Lipschitz equivalent to any homogeneous Carnot group?  

Equivalently, is the mere constancy (up to graded isomorphism) of the nilpotent symbol algebra almost everywhere sufficient to force local metric homogeneity, or can measurable variation in the **norms on \(V_1\)** give rise to distinct global geometries with identical infinitesimal graded structure?

Why is it a "good" problem:  
This question isolates the precise boundary between infinitesimal and global rigidity for sub‑Riemannian and sub‑Finsler geometries.  
Unlike classical rigidity results that assume smooth or analytic constancy of the structure coefficients, this problem asks whether *metric‑measure data alone*—identical tangent Carnot structures up to graded isomorphism—are enough to determine local geometry.  
A positive solution would establish an **infinitesimal rigidity theorem** showing that constancy of the graded symbol algebra controls local geometry, extending the spirit of Cheeger–Keith differentiation rigidity to the non‑commutative Carnot setting.  
A negative solution would uncover the existence of **“measurably homogeneous but geometrically diverse”** sub‑Finsler manifolds, revealing hidden flexibility in metric‑measure structures that share all infinitesimal invariants.  
This formulation captures a simple, fundamental, currently unresolved question connecting geometric control theory, Lie algebra rigidity, and modern analysis on metric spaces—an archetypal “narrow statement, deep reach” kind of problem.